\begin{proposition}[4.8.14]
\label{IV.2.4.8.14-fr}
Sean $k$ un cuerpo, $X$ un $k$-preesquema de tipo finito y $K$ una extensión de
$k$ que contiene una extensión perfecta de $k$. Entonces, el menor cuerpo de
definición del subpreesquema cerrado $(X_{(K)})_{\mathrm{red}}$ de $X_{(K)}$
es una extensión radicial finita de $k$.
\end{proposition}

\begin{proof}
Si $p$ es el exponente característico de $k$, la hipótesis implica que
$k'=k^{p^{-\infty}}$ está contenido en $K$; se sabe que
$(X_{(k')})_{\mathrm{red}}$ es geométricamente reducido sobre $k'$ (4.6.1), por
lo que $(X_{(k')})_{\mathrm{red}}\otimes_{k'}K$ es reducido y, por consiguiente,
igual a $(X_{(K)})_{\mathrm{red}}$ (\textbf I, 5.1.8); en otras palabras, $k'$ es un
cuerpo de definición de $(X_{(K)})_{\mathrm{red}}$. Como $k'$ es una extensión
algebraica de $k$, la conclusión se sigue de (4.8.13).
\end{proof}

\begin{remarks}[4.8.15]
\label{IV.2.4.8.15-fr}
\begin{enumerate}
\item[(i)] La conclusión de (4.8.14) ya no se cumple necesariamente si no se
supone que $K$ contiene una extensión perfecta de $k$. Por ejemplo, sean $k_0$

\oldpage[IV-2]{84}

un cuerpo de característica $p>0$, $k$ el cuerpo de fracciones racionales
$k_0(s,t)$ en $2$ indeterminadas, $L$ el cuerpo
$k(s^{1/p},t^{1/p})$ y $X=\operatorname{Spec}(L)$. Por otra parte, sea $u$
una tercera indeterminada y sea $K$ el cuerpo
$k(u,s^{1/p}+ut^{1/p})$; se comprueba fácilmente que $k$ es algebraicamente
cerrado en $K$ y que $L\otimes_kK$ tiene un nilradical no nulo, por lo que
$X_{(K)}$ no es reducido; pero, como $(X_{(K)})_{\mathrm{red}}$ no está definido
sobre $k$ y $K$ no contiene ninguna subextensión de grado finito sobre $k$
distinta de $k$, el menor cuerpo de definición de $(X_{(K)})_{\mathrm{red}}$ no
puede ser finito sobre $k$.

\item[(ii)] Dado un elemento de uno de los conjuntos de llegada de las
aplicaciones (4.8.2) y una extensión $K_1$ de $K$, son equivalentes las
afirmaciones de que el objeto considerado admite un menor cuerpo de definición
y de que su imagen por $\varepsilon_{K_1,K}$ admite un menor subcuerpo de
definición (siendo necesariamente iguales ambos subcuerpos), como se vio en
(4.8.4).
\end{enumerate}
\end{remarks}
